\chapter{Experimental Protocol}
\label{chap:protocol}

\section{Introduction}

This chapter describes the full experimental setup used across all pruning
experiments in this thesis. It covers the teacher model, the evaluation tasks,
how data is allocated between the search and recovery stages, the hardware
environment, and the reporting conventions. Holding these constant across all
runs ensures that observed differences in accuracy reflect the pruning strategy
alone.

\section{Teacher Model: BERT\textsubscript{BASE}}
\label{sec:teacher}

The teacher is the task-fine-tuned uncased BERT\textsubscript{BASE} model
\citep{devlin2019bert}, built on the Transformer encoder of
\citet{vaswani2017attention}. BERT uses only the encoder side of the
Transformer: there is no decoder, and self-attention is fully bidirectional,
meaning every token attends to every other token in the sequence simultaneously.
This is the property that distinguishes BERT from autoregressive models and that
makes its representations particularly strong for classification tasks.

Each of BERT\textsubscript{BASE}'s twelve layers contains two sublayers. The
first is a multi-head self-attention (MHA) sublayer with twelve heads per layer.
Each head operates on a per-head dimension of $d_k = 64$, obtained by splitting
the hidden dimension $d = 768$ equally across the twelve heads. The second
sublayer is a position-wise feed-forward network (FFN) with an intermediate
width of $d_{\mathrm{ff}} = 3{,}072$, i.e.\ four times the hidden dimension.
Both sublayers use residual connections \citep{he2016deep} and layer
normalisation \citep{ba2016layer}. The full model has approximately 110 million
parameters and uses a WordPiece vocabulary of 30{,}522 uncased tokens
\citep{wu2016googles}.

BERT is pre-trained on BooksCorpus \citep{zhu2015aligning} and English
Wikipedia using two objectives: masked language modelling, which trains the
model to recover randomly masked tokens from their context, and next-sentence
prediction, which trains it to identify whether two sentences are consecutive in
the original document. Fine-tuning then appends a linear classification head on
top of the \texttt{[CLS]} token representation and trains the whole model on the
downstream task for a small number of epochs.

For structured pruning, the relevant architectural units are the attention heads
and the FFN intermediate neurons. Removing a head eliminates its query, key, and
value projections and the corresponding slice of the output projection, yielding
a reduction in multiply--accumulate operations that is predictable from the
mask. Removing a neuron eliminates a row of the first FFN weight matrix and the
corresponding column of the second. Both operations are structured: they produce
dense sub-networks that can be executed efficiently without sparse kernels.

The student is initialised from the same fine-tuned teacher weights and has
the head and neuron masks defined in Chapter~\ref{chap:mask} applied to it.
No additional modules are introduced during search. Recovery then trains the
masked student using the distillation procedure described in
Chapter~\ref{chap:recovery}, with the fine-tuned teacher providing soft targets
and intermediate supervision.

\section{Evaluation Tasks and the GLUE Benchmark}
\label{sec:glue}

\subsection{GLUE}

The General Language Understanding Evaluation benchmark \citep{wang2018glue} is
a collection of nine English natural language understanding tasks spanning
single-sentence classification, pairwise similarity scoring, and natural
language inference. Training set sizes range from roughly 2{,}500 examples
(WNLI) to over 390{,}000 (MNLI), and the linguistic demands range from
surface-level sentiment to fine-grained logical reasoning. The benchmark was
designed to reward broad transfer rather than task-specific overfitting, and it
became the primary evaluation suite for pre-trained language models from BERT
onward.

This work uses two GLUE tasks, selected to provide complementary stress cases
for the pruning method.

\subsection{QNLI}

QNLI is a binary natural language inference task derived from the Stanford
Question Answering Dataset \citep{rajpurkar2016squad}. SQuAD provides Wikipedia
passages paired with crowd-sourced questions, where the answer is a contiguous
span in the passage. GLUE converts this into a classification problem by
splitting each passage into individual sentences and forming (question, sentence)
pairs. A pair is labelled \texttt{entailment} if the sentence contains the
answer, and \texttt{not\_entailment} otherwise. The model must therefore judge,
without access to the rest of the passage, whether a given sentence answers a
given question.

The GLUE release of QNLI contains approximately 104{,}743 training examples and
5{,}463 development examples. The class distribution is approximately balanced
between the two labels. The task is evaluated by accuracy on the development
set.

\subsection{MNLI}

MNLI \citep{williams2018mnli} is a three-way natural language inference task.
Each example consists of a premise and a hypothesis; the label is
\texttt{entailment}, \texttt{neutral}, or \texttt{contradiction}. What
distinguishes MNLI from earlier NLI datasets is its genre diversity: premises
are drawn from ten genres of written and spoken American English, including
fiction, government documents, telephone speech, and travel guides. This
diversity is captured in two separate development sets: a \emph{matched} set,
drawn from the same genres as training, and a \emph{mismatched} set, drawn from
held-out genres.

The training split contains 392{,}702 examples, making it the largest in GLUE.
Each development set contains approximately 9{,}815 examples, with labels
distributed roughly uniformly across the three classes. We report matched
development accuracy as the primary metric; mismatched accuracy is reported
where relevant.

\subsection{Why These Two Tasks}

QNLI and MNLI impose genuinely different demands on a compressed model. QNLI is
binary and smaller, testing whether the pruned model retains the capacity for
cross-sentence semantic matching under a moderate data regime. MNLI is larger
and three-way: discriminating \texttt{neutral} from \texttt{entailment} and
\texttt{contradiction} requires sensitivity to subtle pragmatic and
world-knowledge cues that are among the first capabilities to degrade under
aggressive compression. Together, the two tasks expose failure modes that neither
would reveal alone.

\section{Data Allocation}
\label{sec:data}

The two stages of the pipeline --- search and recovery --- have different data
requirements and use the available data differently.

During search, the goal is to score candidate masks quickly and comparably.
To keep evaluation tractable across thousands of candidates, each mask is
assessed on a small calibration subset. Calibration examples are sampled from
the training split and tokenised to a fixed length; the number of tokens
presented to the surrogate ranges from 2{,}048 to 8{,}192 depending on the
target task and the budget allocated to each search generation. This range was
chosen to balance signal quality against evaluation speed: too few tokens produce
noisy surrogate scores, while too many make the search prohibitively slow.

During recovery, the selected mask is distilled using the full training split
of the corresponding task. Using the complete dataset here is important: the
distillation signal must cover the full input distribution so that the recovered
student does not overfit to the calibration subset seen during search. The
development splits are reserved exclusively for validation and are never used to
guide either search or recovery.

\section{Hardware}
\label{sec:hardware}

All experiments are run on a single consumer laptop equipped with an
NVIDIA RTX 4090 GPU with 16\,GB of VRAM. Mixed precision training is used
throughout the recovery stage where numerically stable \citep{micikevicius2018mixed};
the surrogate metric used during search is computed in full precision, since it
measures small differences in residual updates that would be corrupted by
half-precision rounding.

It is worth noting that running experiments of this scale on a single
consumer GPU is non-trivial. The RTX 4090 is a desktop-class card housed in a
laptop chassis, which means sustained workloads are subject to thermal throttling
that desktop equivalents do not face. Gradient checkpointing and careful batch
sizing were used to keep peak memory within the 16\,GB limit, particularly for
MNLI where long premises can produce large activation tensors.

\section{Baselines and Reporting}
\label{sec:baselines}

All baselines use the same teacher checkpoint, calibration data, target compute
budget, and recovery procedure as the proposed method. The comparison includes
random search, beam search, and published structured pruning results from CoFi
\citep{xia2022cofi} and ZipLM \citep{kurtic2023ziplm}. For each configuration
we report development accuracy, giga-operations at reference sequence length
128, compute keep ratio, number of active heads, number of active neurons, and
number of active layers. Differences within one standard error are treated as
ties. When variance is reported, three seeds are used and the table gives the
mean with standard error; otherwise all runs use a single fixed seed.

\section{Conclusion}
\label{sec:protocol-conclusion}

The protocol described in this chapter provides the controlled foundation on
which the remaining experiments rest. The teacher, the tasks, the data
allocation, the hardware, and the reporting format are fixed. What varies across
the experiments in the next chapter is the search strategy used to find the
pruning mask.

Chapter~\ref{chap:results} presents those results. For each task and each target
compute level, it compares the masks found by different search strategies ---
random, beam, and the proposed evolutionary approach --- after identical recovery.
The protocol established here is what makes those comparisons informative: when
one method outperforms another, there is no alternative explanation on offer.